\section*{第 \ref{ch:g5:water-shapes-the-land} 章——水如何塑造大地}

\begin{solution}{exo:g5:water-shapes-the-land:1}
（a）\omterm{def:g5:water-shapes-the-land:weathering}{风化}：\omterm{def:g3:stones-pebbles-sand:rock}{岩石}在原地被破碎。（b）\omterm{def:g5:water-shapes-the-land:erosion}{侵蚀}：\omterm{def:g3:stones-pebbles-sand:sand}{沙
子}被磨损并带走。（c）\omterm{def:g5:water-shapes-the-land:deposition}{沉积}：泥沉淀下来。（d）
\omterm{def:g5:water-shapes-the-land:weathering}{风化}：树根在原地把石头撑裂。
\end{solution}

\begin{solution}{exo:g5:water-shapes-the-land:2}
\omterm{def:g5:water-shapes-the-land:source-mouth}{源头}是\omterm{def:g1:landscapes-around-us:water}{河流}开始的地方（泉水、融化的冰川
或高处的湿地）。\omterm{def:g5:water-shapes-the-land:source-mouth}{河口}是\omterm{def:g1:landscapes-around-us:water}{河流}终止的地方，位于
大海、\omterm{def:g1:landscapes-around-us:water}{湖泊}或较大的\omterm{def:g1:landscapes-around-us:water}{河流}之中。\omterm{def:g5:water-shapes-the-land:source-mouth}{支流}是较小的\omterm{def:g1:landscapes-around-us:water}{河流}，
汇入较大的\omterm{def:g1:landscapes-around-us:water}{河流}。
\end{solution}

\begin{solution}{exo:g5:water-shapes-the-land:3}
\omterm{def:g5:water-shapes-the-land:source-mouth}{源头}；狭窄的 V 形\omterm{def:g1:landscapes-around-us:landform}{山谷}；\omterm{def:g5:water-shapes-the-land:meander}{曲流}；\omterm{def:g5:water-shapes-the-land:delta}{三角洲}。
\end{solution}

\begin{solution}{exo:g5:water-shapes-the-land:4}
（a）错误：霜、水、树根和冷热交替会缓慢地破碎
即使是坚硬的\omterm{def:g3:stones-pebbles-sand:rock}{岩石}（\omterm{def:g5:water-shapes-the-land:weathering}{风化}）。（b）正确：水流越快，能搬运的碎块就越大。
（c）错误：最重的碎块最先沉降；最细的泥
最后沉降，那时水几乎静止，在\omterm{def:g1:landscapes-around-us:water}{湖泊}或大海中。
\end{solution}

\begin{solution}{exo:g5:water-shapes-the-land:5}
A 位于弯道外侧：那里水流快而深，会冲刷
岸边（\omterm{def:g5:water-shapes-the-land:erosion}{侵蚀}）。B 位于弯道内侧：那里水流
缓慢，并\omterm{def:g5:water-shapes-the-land:deposition}{沉积}下\omterm{def:g3:stones-pebbles-sand:sand}{沙子}（\omterm{def:g5:water-shapes-the-land:deposition}{沉积}）。
\end{solution}

\begin{solution}{exo:g5:water-shapes-the-land:6}
在山区，\omterm{def:g1:landscapes-around-us:water}{河流}陡峭而湍急：它能推动大石块滚动，
并把细小的颗粒带走。靠近大海时，\omterm{def:g1:landscapes-around-us:water}{河流}缓慢：
它已无法搬运大石块，只能搬运沙和泥，
并在那里将它们\omterm{def:g5:water-shapes-the-land:deposition}{沉积}下来。途中，石块也已被破碎、
磨成更小的颗粒。
\end{solution}

\begin{solution}{exo:g5:water-shapes-the-land:7}
（a）上游：$600 \div 6 = 100$ 千米。中游：
$600 \div 2 = 300$ 千米。（b）下游：$600 - 100 - 300 = 200$
千米。由于 $3 \times 200 = 600$，它占\omterm{def:g1:landscapes-around-us:water}{河流}的 $\frac{1}{3}$。
（检验：$\frac{1}{6} + \frac{1}{2} + \frac{1}{3} = \frac{1}{6} +
\frac{3}{6} + \frac{2}{6} = \frac{6}{6}$。）
\end{solution}

\begin{solution}{exo:g5:water-shapes-the-land:8}
浅色斑块是\omterm{def:g5:water-shapes-the-land:deposition}{沉积}在弯道内侧的沙和\omterm{def:g3:stones-pebbles-sand:pebble}{砾石}，
那里水流缓慢（\omterm{def:g5:water-shapes-the-land:deposition}{沉积}）。深色弯曲的\omterm{def:g1:landscapes-around-us:water}{湖泊}是旧河湾，
即两个河湾逐渐靠近，\omterm{def:g1:landscapes-around-us:water}{河流}冲破了它们之间的陆地颈部，
走了捷径，将旧河湾留下，成为
\omterm{def:g1:landscapes-around-us:water}{湖泊}。
\end{solution}

\begin{solution}{exo:g5:water-shapes-the-land:9}
河边平坦的草地是\omterm{def:g5:water-shapes-the-land:flood}{河漫滩}，即\omterm{def:g1:landscapes-around-us:water}{河流}泛滥时所淹没的
土地。它们之所以平坦，正是因为\omterm{def:g1:landscapes-around-us:water}{河流}用自身的\omterm{def:g3:stones-pebbles-sand:sand}{泥沙}堆积而成。
总有一天\omterm{def:g1:landscapes-around-us:water}{河流}会再次泛滥，把学校淹没，
而湍急的\omterm{def:g5:water-shapes-the-land:flood}{洪水}是危险的。学校应建在较高的地方，
草地则保留为农田、草地或公园。
\end{solution}

\begin{solution}{exo:g5:water-shapes-the-land:10}
（a）一年有 12 个月：$120$ 亿 $\div\, 12 = 10$ 亿吨，
即平均每月 10 亿吨。（b）10 亿是 $1\,000\,000\,000$；除以 10，
得 $100\,000\,000$，即一亿次卡车的运载量。
\end{solution}

\begin{solution}{exo:g5:water-shapes-the-land:11}
一种可能的故事。这颗颗粒原是高山上\omterm{def:g3:stones-pebbles-sand:rock}{岩石}的一部分。霜
冻一夜又一夜地在它周围撑开裂缝，直到一块\omterm{def:g3:stones-pebbles-sand:rock}{岩石}脱落：这就是\omterm{def:g5:water-shapes-the-land:weathering}{风化}。
碎块落入下方的溪流。在上游，
水流湍急而陡峭：它把碎块推着滚动（\omterm{def:g5:water-shapes-the-land:erosion}{侵蚀}），
并在与其他石头相撞时，碎成\omterm{def:g3:stones-pebbles-sand:pebble}{卵石}，再碎成\omterm{def:g3:stones-pebbles-sand:sand}{沙粒}。
在中游，\omterm{def:g1:landscapes-around-us:water}{河流}变缓；颗粒在\omterm{def:g5:water-shapes-the-land:meander}{曲流}弯道内侧的沙滩上停留一段时间，
直到\omterm{def:g5:water-shapes-the-land:flood}{洪水}再次将它卷起。在
下游，\omterm{def:g1:landscapes-around-us:water}{河流}缓慢而宽阔；颗粒随水流一同抵达
\omterm{def:g5:water-shapes-the-land:source-mouth}{河口}。那里水流几乎停止，颗粒便
与\omterm{def:g5:water-shapes-the-land:delta}{三角洲}的泥和沙一起沉降：这就是\omterm{def:g5:water-shapes-the-land:deposition}{沉积}。
\end{solution}
